
\subsubsection{Key Findings Interpretation}

\textbf{Methodology Works for Simulated Signature Detection}: Simulated Cohen's d=7.835 indicates the statistical framework can reliably detect alignment method patterns when they exist at large magnitudes in designed data. Perfect simulated alignment purity (1.000) demonstrates the clustering pipeline functions correctly. However, this validates methodology capability (can the framework detect patterns?), not empirical reality (do real models exhibit such patterns?).

\textbf{Framework Can Profile Real Models}: The minimal real-model validation (10 tasks, 300 generations) confirms the framework can extract actual performance signatures from real models. The execution-aligned phi-2 achieved measurably higher correctness (13.0\%) than smaller models (1.0\%, 0.0\%). However, the $8\times$ scale difference prevents attributing this to alignment rather than capacity.

\textbf{Execution-Dominance Pattern Detection}: The 0.0\% percentile rank on correctness for phi-2 confirms the framework can identify correctness-specialization. However, the model scale confound means even this successful result may partially reflect capacity differences. A cleaner validation comparing execution vs preference models at matched scale would isolate pattern recognition capability from confounding factors.

\textbf{Preference-Balance Detection Remains Unvalidated}: The M2 failure (53.3\% > 30\% threshold) cannot distinguish between ''framework cannot detect balanced patterns,'' ''preference pattern poorly designed,'' or ''model scale confound overwhelms balanced pattern.'' Until matched-scale validation, M2 capability remains an open question.

\subsubsection{Limitations and Future Work}

\#\#\#\# POC Validation with Simulated Data (Critical Limitation)

The primary proof-of-concept validation used simulated performance data to demonstrate methodology feasibility, not real model inference at scale. Simulated data was designed with differentiated signatures to test whether clustering analysis detects these patterns.

\textbf{Why this matters}: Simulated Cohen's d=7.835 validates that the statistical framework works (can detect signatures when they exist), but does not validate that real models exhibit signatures at all, let alone at this magnitude. Real model inference at scale may yield smaller effect sizes, noisier clustering, different signature patterns, or no detectable signatures.

\textbf{Mitigation path}: Validate with full real-model inference (164 tasks $\times$ 30 samples $\times$ 4+ models). If real Cohen's d > 1.5, existence claim is validated. If real d < 1.5, existence hypothesis is refuted.

\textbf{Minimal real validation conducted}: The 10-task, 3-model validation (300 generations) confirms framework functionality but lacks statistical power for robust conclusions. This validates that the pipeline works on real data but does not test the hypothesis at scale.

\#\#\#\# Model Scale Confound (Affects Both M1 and M2)

Comparing 2.7B execution model (phi-2) against 350M preference-categorized model (codegen-350M-mono) conflates alignment method with model capacity. Prior work shows model scale strongly affects performance.

\textbf{Impact on M2}: Cannot distinguish ''methodology cannot detect balanced patterns'' from ''small-scale model performs poorly.'' M2 failure is uninterpretable without matched-scale validation.

\textbf{Impact on M1}: Even the successful M1 result may partially reflect scale confound. The execution model's correctness dominance could stem from $8\times$ capacity advantage rather than pure correctness-specialization.

\textbf{Mitigation path}: Compare alignment methods at matched scales (2B execution vs 2B preference, 350M execution vs 350M preference, 7B execution vs 7B preference).

\#\#\#\# Small Sample Size

Testing 3 total models (vs. planned 6-8) reduces statistical power substantially. With N=3, confidence intervals are wide and edge cases may not generalize. Scaling to 10+ models per category is required for robust claims.

\#\#\#\# Single Benchmark

HumanEval+ is standardized but function-level. Repository-level benchmarks (BigCodeBench), real-world tasks (NaturalCodeBench), or domain-specific evaluations may reveal different signature patterns. Cross-benchmark validation remains untested.

\#\#\#\# Temperature Confound

Results used T=0.8 for generation but did not test signature stability across temperature settings. High temperature increases diversity, potentially inflating signature detectability. Signature stability across T $\in$ [0.2, 1.0] is untested.

\subsubsection{Broader Impact}

\textbf{Positive}: If validated with real models at scale, the signature detection framework could enable practitioners to diagnose model optimization biases before deployment. Signature profiling could identify whether a candidate model has undesirable specialization. This would support informed method selection beyond leaderboard rankings.

\textbf{Enabling research}: Signature-based benchmark validation becomes possible if the methodology validates with real models. If a benchmark claims to measure ''code quality,'' but models trained for correctness alone dominate the benchmark, the benchmark may not measure intended dimensions.

\textbf{Potential Negative}: Revealing optimization biases in deployed models may require costly retraining or model replacement. However, transparency about biases is likely net-positive for responsible AI deployment.

\subsubsection{Conceptual Contributions}

Beyond proof-of-concept validation, this work introduces:

\textbf{1. Alignment Methods as Implicit Objectives}: Reframing alignment not as ''training procedures'' but as ''implicit objective function definitions'' enables new research questions.

\textbf{2. Signature-Based Method Selection}: If validated with real models, signature profiling could enable multi-criteria selection beyond single-metric leaderboards.

\textbf{3. Evaluation Ontology}: The framework positions signature detection as infrastructure for mapping (alignment method $\times$ objective dimension $\times$ benchmark sensitivity).

\subsubsection{Future Directions}

\textbf{Immediate: Real-Model Validation at Scale}: The most critical next step is validating the methodology with full model inference on HumanEval+ (164 tasks, 30 samples, 4+ models). This would determine whether real models exhibit detectable signatures and at what effect sizes.

\textbf{Matched-Scale Comparisons}: Testing alignment methods at matched scales is essential to isolate alignment effects from capacity effects.

\textbf{Cross-Benchmark Validation}: Extending to MBPP+, BigCodeBench, LiveCodeBench would test signature stability.

\textbf{Temperature Sensitivity}: Testing across temperature settings (T $\in$ [0.2, 1.0]) would validate that signatures reflect model properties rather than sampling artifacts.

\textbf{Signature-Guided Alignment}: If real-model validation succeeds, can we design alignment to produce target signatures?

\subsubsection{Honest Assessment}

This work proposes a diagnostic framework for alignment signature detection and validates through proof-of-concept that the methodology can successfully identify patterns when they exist in simulated data. A minimal real-model validation confirms the framework can profile actual models. However, critical limitations---proof-of-concept simulation without full-scale real-model validation, model scale mismatch affecting both M1 and M2, small samples, temperature confound---prevent definitive claims about real alignment method behavior.

This represents an important first step, not a complete validation. Methodology proof-of-concept (demonstrating signature detection is technically feasible) precedes full validation (demonstrating signatures exist robustly in real models). The immediate real-model validation experiment at scale will either strengthen claims (signatures persist in reality) or refute them (signatures are weaker than simulated data suggests or absent entirely).

\textbf{What we can claim}: We developed a signature detection framework and demonstrated through proof-of-concept that it can identify alignment method patterns when they exist in data at large effect sizes. Minimal real-model validation confirms the framework can profile actual models.

\textbf{What we cannot claim}: That real models exhibit signatures at scale, that signatures exist at any particular magnitude in practice, that the framework can detect all signature types (balanced patterns unvalidated), or that results generalize across scales/temperatures/benchmarks.

\textbf{Path forward}: Full-scale real-model validation (164 tasks, matched scales, 10+ models) determines whether this methodology transitions from ''proposed framework'' to ''validated diagnostic tool.''
